% !TITLE 红豆曲
% !AUTHOR 曹雪芹
% !DYNASTY 明清
% !GENRE 曲
% !ORDER 8
% !SOURCE 《红楼梦》（借贾宝玉之手）
% !EDITION 《红楼梦》，中国艺术研究院红楼梦研究所校注，人民文学出版社（2008年版为候选）
% !EDITION_STATUS pending

\begin{poem}
滴不尽相思血泪抛红豆\textsuperscript{1}，开不完春柳春花满画楼，
睡不稳纱窗风雨黄昏后，忘不了新愁与旧愁，
咽不下玉粒金莼\textsuperscript{2}噎满喉\textsuperscript{3}，照不见菱花镜\textsuperscript{4}里形容\textsuperscript{5}瘦。
展不开的眉头，捱不明\textsuperscript{6}的更漏\textsuperscript{7}。
呀！恰便似遮不住的青山隐隐，流不断的绿水悠悠。
\end{poem}

\begin{annotations}
\notetext{1}{红豆：红豆树（相思木）的种子，色红，古代常用作爱情信物，又名“相思子”。王维《相思》：“红豆生南国，春来发几枝。愿君多采撷，此物最相思。”}
\notetext{2}{玉粒金莼：玉粒指白米（像玉一样的米粒）；金莼（chún）指莼菜（一种水生蔬菜，江南名产）。代指精美的饭菜。}
\notetext{3}{噎满喉：堵在喉咙里咽不下去。因相思而茶饭不思。}
\notetext{4}{菱花镜：古代铜镜背面常刻菱花图案，故称“菱花镜”。}
\notetext{5}{形容：容貌体态。镜子里照见的自己日渐消瘦。}
\notetext{6}{捱不明：熬不到天亮。捱（ái），忍受、苦熬。}
\notetext{7}{更漏：古代漏壶计时，一夜分五更，故称“更漏”。听着一更一更的更漏声睡不着，天总也不亮。}
\end{annotations}

\begin{appreciation}
这支曲子出自《红楼梦》第二十八回，是全书写相思最浓的笔墨之一。冯紫英设宴，宝玉、薛蟠、蒋玉菡一干人在席上行酒令，轮到宝玉，他开口唱了这支后世题为《红豆曲》的小调——满座酒酣耳热，他唱出来的却是满纸血泪。

曲名落在红豆二字上头，首句滴不尽相思血泪抛红豆正把它含在句中。红豆又名相思子，色赤而坚，古人拿它当定情的信物；王维那首《相思》我们读过，红豆生南国，春来发几枝，愿君多采撷，此物最相思。泪本是清的，这里偏偏写成血泪，又把它和红豆并在一处，一颗一颗落下来的分不清是泪还是豆，形、色、情三样都在里头了。

往下是一串不字：开不完春柳春花满画楼，睡不稳纱窗风雨黄昏后，忘不了新愁与旧愁，咽不下玉粒金莼噎满喉，照不见菱花镜里形容瘦。花是满楼开着的她偏偏看不下去，玉粒金莼是上等饭食她却咽不下去，镜子本是照人的她一样不敢去照。相思到底连吃饭照镜这两样最寻常的事都做不成。

展不开的眉头，捱不明的更漏——句子忽然收短，一更一更地熬，天偏偏不亮。这个瘦字，《行行重行行》里汉代那位思妇早已用过，相去日已远，衣带日已缓。一个在汉末一个在大观园，相隔一千多年，为相思瘦下去的模样原来古今分毫不差。

末一句是全曲的顶：呀！恰便似遮不住的青山隐隐，流不断的绿水悠悠。一个呀字露了底——酒席上唱曲的人，唱着唱着忘了自己是在唱曲。青山遮不住绿水流不断这两句原也有来历，杜牧写过青山隐隐水迢迢，原是写扬州的，曹雪芹借它来收相思，恰好相思也是遮不住流不断的。

宝玉是书中最懂得体贴女孩心事的人，这支曲偏偏安排在众声喧哗的宴席上唱出来。《红楼梦》惯于把最悲的调子压进最热闹的场面里，读的时候不觉，回头再看，满席的灯火倒像替一个人守了一整夜。
\end{appreciation}
